\section{Preliminaries}
\label{sec:prelim}

This section fixes the two mathematical tools the rest of the paper
builds on: uniform B-splines as a continuous-time trajectory
representation, both in $\mathbb{R}^3$ and, in cumulative form, on
$\SO$; and the \ac{map} estimator that fits them to the measurements as
a weighted nonlinear least-squares problem.  A reader already fluent in
basalt-style splines and factor-graph \ac{map} estimation can skip to
\autoref{sec:overview}.

\subsection{Uniform B-splines in $\mathbb{R}^3$}
\label{sec:prelim_bspline}
An order-$N$ (degree $N{-}1$) uniform B-spline with control points
$\{\mathbf{P}_i\}\subset\mathbb{R}^3$ and constant knot spacing $\Delta t$
represents a curve as a weighted sum of basis functions,
\begin{equation}
  \vp(t) = \sum_i \mathbf{P}_i\, B_{i,N}(t),
  \label{eq:bspline_def}
\end{equation}
where the $B_{i,N}$ are the Cox--de Boor basis
functions~\cite{furgale2013unified}.  Two properties make this
representation attractive for asynchronous sensor fusion.

\emph{Local support.}  At any instant only $N$ consecutive control
points have nonzero basis, so on the segment $t\in[t_i,t_{i+1})$, with
normalized time $u=(t-t_i)/\Delta t\in[0,1)$, the curve is a fixed
blending of $N$ control points,
\begin{equation}
  \begin{aligned}
  \vp(t) &= \sum_{j=0}^{N-1} B_j(u)\,\mathbf{P}_{i+j}, \\
  \bigl[B_0(u)\,\cdots\,B_{N-1}(u)\bigr]
     &= \mathbf{u}^{\top}\mathbf{M}^{(N)},
  \end{aligned}
  \label{eq:bspline_blend}
\end{equation}
with $\mathbf{u}=[1,u,\dots,u^{N-1}]^{\top}$ and a constant $N\times N$
blending matrix $\mathbf{M}^{(N)}$~\cite{usenko2020basalt}.  Local
support is what bounds each measurement's Jacobian to $N$ knots and is
the property the marginalization rule of \autoref{sec:sw}
(Proposition~\ref{prop:marg}) exploits.

\emph{Analytic derivatives.}  Differentiating the basis in
\eqref{eq:bspline_blend} yields the velocity $\dot{\vp}(t)$ and
acceleration $\ddot{\vp}(t)$ in closed form at \emph{any} $t$, as
lower-order splines over the same control points, with no finite
differencing or inter-sample interpolation.  This is what lets the
accelerometer factor (which needs $\ddot{\vp}$) and the min-snap
regularizer (which needs $\vp^{(4)}$) be evaluated exactly at each
measurement's timestamp.  We use a quintic ($N{=}6$) position spline so
that snap is well defined and continuous.

\subsection{Cumulative $\SO$ B-splines}
\label{sec:prelim_so3}
Orientation lives on the manifold $\SO$, on which a weighted sum of
rotations $\sum_i \vR_i B_i$ is not itself a rotation, so
\eqref{eq:bspline_def} does not transfer directly.  The
\emph{cumulative} parameterization of Sommer et
al.~\cite{sommer2020efficient}, used in the basalt
\ac{vio}~\cite{usenko2020basalt}, instead writes the orientation as a
base rotation followed by a product of incremental rotations, each
mapped back onto the manifold by the exponential map,
\begin{equation}
  \begin{aligned}
  \vR(t) &= \vR_{i}
           \prod_{j=1}^{N-1}
           \Exp\!\bigl(\tilde{B}_j(u)\,\vOm_j\bigr), \\
  \vOm_j &= \Log\!\bigl(\vR_{i+j-1}^{-1}\vR_{i+j}\bigr)\in\so,
  \end{aligned}
  \label{eq:prelim_so3}
\end{equation}
where the increments $\vOm_j$ are the relative rotations between
consecutive \emph{absolute} knots $\vR_i\in\SO$ (the optimization
variables, basalt convention), and the cumulative basis is the tail sum
of the ordinary basis of \eqref{eq:bspline_blend},
\begin{equation}
  \tilde{B}_j(u) = \sum_{l=j}^{N-1} B_l(u).
  \label{eq:cumbasis}
\end{equation}
Because each $\Exp(\cdot)$ returns a rotation, $\vR(t)\in\SO$ holds by
construction for any knots.  This form keeps the intuitive absolute
control mesh, preserves the $N$-knot local support of the $\mathbb{R}^3$
case exactly (again used by Proposition~\ref{prop:marg}), and yields the
body angular velocity $\vom_b(t)$ in closed form by differentiating the
product (basalt's \texttt{evaluate\_lie()}), with no numerical
differentiation.  We use a cubic ($N{=}4$) orientation spline;
\eqref{eq:ori_spline} in \autoref{sec:state} is this construction with
$N{=}4$ (base knot $\vR_{k-3}$, three increments).

\subsection{\ac{map} Estimation as Weighted Least Squares}
\label{sec:prelim_map}
The control points and biases above form the finite-dimensional state
$\calX$ (\autoref{sec:problem}).  Each measurement $\vz_i$ relates to
the state through a known model $\mathbf{h}_i$,
$\vz_i = \mathbf{h}_i(\calX) + \mathbf{n}_i$, with the noise
$\mathbf{n}_i$ zero-mean Gaussian of covariance
$\boldsymbol{\Sigma}_i$ and independent across measurements.  The
likelihood is then
$p(\vz\mid\calX)\propto\prod_i
 \exp(-\tfrac12\mathbf{r}_i^{\top}\boldsymbol{\Sigma}_i^{-1}\mathbf{r}_i)$
with residuals $\mathbf{r}_i=\vz_i-\mathbf{h}_i(\calX)$, so its negative
logarithm is a weighted sum of squared residuals: maximum-likelihood
estimation minimizes
$\tfrac12\sum_i\mathbf{r}_i^{\top}\boldsymbol{\Sigma}_i^{-1}\mathbf{r}_i$,
in which each residual is weighted by its inverse covariance
(information) $\boldsymbol{\Sigma}_i^{-1}$.  For a scalar residual this
is a scalar weight $\lambda_i=1/\sigma_i^2$: a more certain sensor is
trusted more.  The per-stream weights $\lambda_a,\lambda_g,\lambda_\psi$
of \autoref{sec:models}, and the rate-adaptive law that sets them, are
precisely these information weights.

Prior knowledge $p(\calX)$, here trajectory smoothness and, online, the
marginalization prior summarizing past data, turns the estimate into a
\ac{map} one, $\arg\max_{\calX} p(\vz\mid\calX)\,p(\calX)$; its negative
log contributes the extra quadratic terms $E_{\text{reg}}$ and
$E_{\text{prior}}$.  Radar Doppler outliers are heavy-tailed and break
the Gaussian assumption, so the quadratic on the Doppler residual is
replaced by a Huber kernel $\rho_\delta$ (quadratic for
$|\mathbf{r}|<\delta$, linear beyond), which bounds each outlier's
influence while staying Gaussian-optimal on the inliers, a robust
M-estimator.  Collecting all terms gives the weighted nonlinear
least-squares objective instantiated for our sensors in
\eqref{eq:map}.  We minimize it with Levenberg--Marquardt: each
iteration linearizes the residuals about the current estimate,
$\mathbf{r}_i(\calX\boxplus\boldsymbol{\delta})\approx\mathbf{r}_i
 +\mathbf{J}_i\boldsymbol{\delta}$, and solves the damped normal
equations
$\bigl(\sum_i\mathbf{J}_i^{\top}\mathbf{W}_i\mathbf{J}_i
 +\mu\mathbf{I}\bigr)\boldsymbol{\delta}
 =-\sum_i\mathbf{J}_i^{\top}\mathbf{W}_i\mathbf{r}_i$
for the on-manifold update $\boldsymbol{\delta}$, with $\mathbf{W}_i$
the (Huber-reweighted) information and $\mu$ the damping.
Here $\boxplus$ denotes the generalized addition that applies
a small update to each part of the state in its own geometry: ordinary
addition for the vector components, and $\vR\mapsto\vR\,\Exp(\boldsymbol{\delta})$
for the rotations, so the updated estimate stays a valid rotation.  We supply
analytic Jacobians $\mathbf{J}_i$.  The continuous-time representation
closes the loop here: every $\mathbf{J}_i$ needs the spline value and
its derivatives at the measurement's exact timestamp, all available in
closed form from \autoref{sec:prelim_bspline} and
\autoref{sec:prelim_so3}.
